% !PRINTED 54
\weberpage{54}

3. Ist $u_0$ ein Unstetigkeitspunkt, so gibt es eine \gesperrt{ganze positive} Zahl $m$ von der Art, daß $(u-u_0)^m\varphi(u)$ in dem Punkte $u_0$ endlich, stetig und von Null verschieden bleibt und daher nach ganzen positiven Potenzen von $u - u_0$ entwickelbar ist. Der Unstetigkeitspunkt wird in diesem Falle \gesperrt{von der $m$ten Ordnung} genannt. Es ergibt sich daraus eine Entwickelung von $\varphi(u)$ nach steigenden Potenzen von $u - u_0$, die mit $(u-u_0)^{-m}$ anfängt und in einem um $u_0$ beschriebenen Kreise, der bis zum nächstgelegenen Unstetigkeitspunkte reicht, konvergiert. Der Koeffizient von $(u-u_0)^{-1}$ in dieser Entwickelung heißt das \gesperrt{Residuum} dieser Funktion für den Punkt $u_0$.

4. \gesperrt{Der Cauchysche Satz.} Das Integral
\[
\frac{1}{2\pi i}\int \varphi(u)\,du,
\]
in positivem Sinne über die Begrenzung eines endlichen Flächenstückes erstreckt, das auf der Randlinie keine Unstetigkeitspunkte enthält, ist gleich der Summe der Residuen für die im Inneren des Flächenstückes liegenden Unstetigkeitspunkte. Als positive Integrationsrichtung ist dabei diejenige anzusehen, die gegen das Innere des Flächenstückes so liegt, wie die $v$-Achse zur $w$-Achse, so daß bei der üblichen Bestimmungsweise dieser Achsen beim positiven Durchlaufen des Randes das Innere der Fläche zur Linken bleibt.

5. Die Funktion
\[
\psi(u) = \frac{d\log\varphi(u)}{du} = \frac{1}{\varphi(u)}\frac{d\varphi(u)}{du}
\]
hat nur Unstetigkeitspunkte erster Ordnung, und wenn in einem Punkte $u_0$ das Produkt $(u-u_0)^{-m}\varphi(u)$ endlich und von Null verschieden ist, so ist $m$ das Residuum von $\psi(u)$ für diesen Punkt. Ist $m$ positiv, so heißt $u_0$ ein Nullpunkt $m$ter Ordnung von $\varphi(u)$.

Hiernach ist eine unmittelbare Folgerung des \gesperrt{Cauchyschen} Theorems:

6. Das Integral
\[
\frac{1}{2\pi i}\int d\log\varphi(u),
\]
ausgedehnt in positiver Richtung über die Begrenzung eines Flächenstückes, ist gleich der Anzahl der Nullpunkte, vermindert um die Anzahl der Unstetigkeitspunkte, die im Inneren dieses
%CONT>>
